% Part: many-valued-logic
% Chapter: infinite-valued-logics
% Section: introduction

\documentclass[../../../include/open-logic-section]{subfiles}

\begin{document}

\olfileid{mvl}{inf}{int}

\olsection{مقدمة}

ليس من شرط المصفوفة أن تكون قيم الصدق فيها متناهية. ومن أظهر
الخيارات لمجموعة لا متناهية مجموعة الأعداد النسبية المحصورة بين
${\OLMathNumeral{0}}$ و${\OLMathNumeral{1}}$، وهي~${\OLId{latin-looped}{V}}_\infty = [{\OLMathNumeral{0}},{\OLMathNumeral{1}}] \cap \Rat$؛ أي
\begin{align*}
    {\OLId{latin-looped}{V}}_\infty & = \Setabs{\frac{{\OLId{latin-ordinary}{n}}}{{\OLId{latin-ordinary}{m}}}}{{\OLId{latin-ordinary}{n}},{\OLId{latin-ordinary}{m}} \in \Nat \text{ و } {\OLId{latin-ordinary}{m}}>{\OLMathNumeral{0}} \text{ و } {\OLId{latin-ordinary}{n}}\le {\OLId{latin-ordinary}{m}}}.
\intertext{وعند النظر في مجموعة قيم الصدق غير المنتهية هذه، يفيد غالبًا
النظر أيضًا، لكل عدد طبيعي~${\OLId{latin-ordinary}{m}}\ge{\OLMathNumeral{2}}$، في المجموعات الجزئية}
{\OLId{latin-looped}{V}}_{\OLId{latin-ordinary}{m}} & = \Setabs{\frac{{\OLId{latin-ordinary}{n}}}{{\OLId{latin-ordinary}{m}}-{\OLMathNumeral{1}}}}{{\OLId{latin-ordinary}{n}} \in \Nat \text{ و } {\OLId{latin-ordinary}{n}}<{\OLId{latin-ordinary}{m}}}
\intertext{فمثلًا، ${\OLId{latin-looped}{V}}_{\OLMathNumeral{5}}$ هي المجموعة ذات قيم الصدق الخمس المتساوية التباعد،}
{\OLId{latin-looped}{V}}_{\OLMathNumeral{5}} & = \{{\OLMathNumeral{0}}, \frac{{\OLMathNumeral{1}}}{{\OLMathNumeral{4}}}, \frac{{\OLMathNumeral{1}}}{{\OLMathNumeral{2}}}, \frac{{\OLMathNumeral{3}}}{{\OLMathNumeral{4}}}, {\OLMathNumeral{1}}\}.
\end{align*}
وفي المنطقات المبنية على هذه المجموعات لا تكون، في العادة، إلا~${\OLMathNumeral{1}}$
مميزة؛ أي~${\OLId{latin-looped}{V}}^+ = \{{\OLMathNumeral{1}}\}$. فـ~${\OLMathNumeral{1}}$ تقوم مقام الصدق المطلق، و${\OLMathNumeral{0}}$ مقام
الكذب المطلق، ويجوز مع ذلك لـ~!!{formula}s أن تأخذ أي قيمة متوسطة
في~${\OLId{latin-looped}{V}}$.

ويجوز بدل النسبية أن نأخذ~${\OLId{latin-looped}{V}}_{[{\OLMathNumeral{0}},{\OLMathNumeral{1}}]} = [{\OLMathNumeral{0}},{\OLMathNumeral{1}}]$، وهي مجموعة الأعداد
\emph{الحقيقية} بين~${\OLMathNumeral{0}}$ و${\OLMathNumeral{1}}$، أو أن نأخذ مجموعة جزئية أخرى لا
متناهية من~$[{\OLMathNumeral{0}},{\OLMathNumeral{1}}]$. والمنطقات التي تتخذ مثل هذه المجموعات قيمًا
للصدق تسمى غالبًا \emph{ضبابية}.


\end{document}

